\section{函数列与函数项级数}

\subsection{函数列的一般概念}

\begin{definition}[函数列的收敛]
    设 $\{S_n(x)\}$ 是 $I$ 上的函数序列，若 $\forall\,x_0 \in I$，都有 $\{S_n(x_0)\}$ 收敛于 $S(x_0)$，那么称 $S_n(x)$ 在区间 $I$ 上\textbf{收敛}或\textbf{逐点收敛}。
\end{definition}

\begin{definition}[函数列的一致收敛]
    设 $\{S_n(x)\}$ 是 $I$ 上的函数序列，若
    $$
    \forall\,\varepsilon > 0: \exists\,N \in \mathbb{N}^{+}: \forall\,n \ge N,\ x \in I: |S_n(x) - S(x)| < \varepsilon
    $$
    那么称 $\{S_n(x)\}$ \textbf{一致收敛}于 $S(x)$。
\end{definition}

下面给出函数列一致收敛的一个充要条件。

\begin{theorem}
    设 $\{S_n(x)\}$ 是 $I$ 上的函数列，$S(x)$ 是 $I$ 上的函数。定义数列
    $$
    \beta_n = \sup_{x \in I} \abs{S_n(x) - S(x)}
    $$
    那么，函数列 $\{S_n(x)\}$ 在 $I$ 上一致收敛的充要条件是 $\lim\limits_{n \to \infty} \beta_n = 0$。

    \begin{proof}
        \begin{itemize}
            \item 必要性：设 $\{S_n(x)\}$ 一致收敛与 $S(x)$，那么根据定义：
            $$
            \forall\,\varepsilon > 0: \exists\,N \in \mathbb{N}^{+}: \forall\,n \ge N,\ x \in I: |S_n(x) - S(x)| < \frac{\varepsilon}{2}
            $$
            而 $\beta_n = \sup\limits_{x \in I} \abs{S_n(x) - S(x)} \le \frac{\varepsilon}{2} < \varepsilon$，所以 $\lim\limits_{n \to \infty} \beta_n = 0$。

            \item 充分性：设 $\lim\limits_{n \to \infty} \beta_n = 0$，那么根据定义：
            $$
            \forall\,\varepsilon > 0: \exists\,N \in \mathbb{N}^{+}: \forall\,n \ge N: \beta_n < \varepsilon
            $$
            根据定义 $\abs{S_n(x) - S(x)} \le \sup \abs{S_n(x) - S(x)} = \beta_n$，所以 $\{S_n(x)\}$ 在 $I$ 上一致收敛于 $S(x)$。
        \end{itemize}
    \end{proof}
\end{theorem}

同样，函数列也有 Cauchy 收敛原理。

\begin{theorem}[函数列的 Cauchy 收敛原理]
    函数列 $\{S_n(x)\}$ 在区间 $I$ 上一致收敛的充要条件是
    $$
    \forall\,\varepsilon > 0: \exists\,N \in \mathbb{N}^{+}: \forall\,n, m \ge N,\ x \in I: \abs{S_n(x) - S_m(x)} < \varepsilon
    $$
\end{theorem}

\subsection{作业}

\begin{problem}
    习题 11.3.1
    \begin{solution}
        \begin{itemize}
            \item[\textbf{2)}] 对于 $n, p \in \mathbb{N}^{+}$，有
            $$
            \begin{aligned}
                \sum_{k = n + 1}^{n + p} \abs{\frac{\sin k}{2^{k}}} \le \sum_{k = n + 1}^{n = p} \frac{1}{2^{k}} < \frac{1}{2^{n}}
            \end{aligned}
            $$
            当 $n \to \infty$ 时该值趋近于 $0$，因此原级数绝对收敛。

            \item[\textbf{3)}] 对于 $n, p \in \mathbb{N}^{+}$，有
            $$
            \begin{aligned}
                \sum_{k = n + 1}^{n + p} \abs{\frac{\cos n + \sin n}{n (n + \sin n)}} & = \sum_{k = n + 1}^{n + p} \abs{\frac{\sqrt{2} \sin \ab(n + \frac{\pi}{4})}{n (n + \sin n)}} \\
                & \le \sum_{k = n + 1}^{n + p} \frac{\sqrt{2}}{n (n - 1)} \\
                & = \frac{\sqrt{2}}{n} - \frac{\sqrt{2}}{n + p} < \frac{\sqrt{2}}{n}
            \end{aligned}
            $$
            当 $n \to \infty$ 时该值趋近于 $0$，因此原级数绝对收敛。
        \end{itemize}
    \end{solution}
\end{problem}

\begin{problem}
    习题 11.3.4
    \begin{solution}
        根据 Taylor 公式，有
        $$
        f(x) = f(0) + f'(0) x + O(x^2)
        $$
        根据 $\lim\limits_{n \to 0} \frac{f(x)}{x} = 0$ 得到 $f(0) = f'(0) = 0$。因此 $f(x) = O(x^2)$。

        于是 $\sum\limits_{n = 1}^{\infty} \abs{f(\frac{1}{n})} = \sum\limits_{n = 1}^{\infty} O\ab(\frac{1}{n^2})$。因此该级数绝对收敛。
    \end{solution}
\end{problem}

\begin{problem}
    习题 11.3.5
    \begin{solution}
        \begin{enumerate}
            \item[\textbf{2)}] 设 $f(x) = \frac{\ln^2 x}{x}$，那么 $f'(x) = \frac{2 \ln x - \ln^2 x}{x^2}$。当 $x > \ln 2$ 时 $f'(x) < 0$，因此 $\ab\{\frac{\ln^2 n}{n}\}$ 是一个递减数列。同时显然有 $\lim\limits_{n \to \infty} \frac{\ln^2 n}{n} = 0$，因此根据 Lebniz 判别法，$\sum\limits_{n = 1}^{\infty} (-1)^{n} \frac{\ln^2 n}{n}$ 收敛。而 $\abs{(-1)^{n} \frac{\ln^2 n}{n}} = O\ab(\frac{1}{n})$，因此原级数条件收敛。

            \item[\textbf{4)}] $\lim\limits_{n \to \infty} (-1)^{n - 1} \frac{1}{\sqrt[n]{n}}$ 发散，因此级数发散。
            
            \item[\textbf{6)}] $\{\E - \ab(1 + \frac{1}{n})^{n}\}$ 是一个递减数列，并且 $\lim\limits_{n \to \infty} \ab(\E - \ab(1 + \frac{1}{n})^{n}) = 0$，根据 Lebniz 判别法原级数收敛。
        \end{enumerate}
    \end{solution}
\end{problem}

\begin{problem}
    习题 11.3.7
    \begin{proof}
        设 $S_n$ 为 $y_n$ 的部分和，那么由题知 $\lim\limits_{n \to \infty} S_n$ 收敛。于是：
        $$
        \begin{aligned}
            \sum_{k = 1}^{n} x_n y_n & = \sum_{k = 1}^{n} x_n (S_n - S_{n - 1}) \\
            & = x_n S_n + \sum_{k = 1}^{n - 1} S_{k} (x_k - x_{k + 1})
        \end{aligned}
        $$

        由 $\sum\limits_{n = 1}^{\infty} (x_n - x_{n + 1})$ 绝对收敛可知 $\lim\limits_{n \to \infty} x_n$ 收敛，因此 $\lim\limits_{n \to \infty} x_n S_n$ 收敛。

        此外，因为 $S_n$ 收敛，其必定有界，因此
        $$
        \sum_{k = 1}^{n - 1} \abs{S_{k} (x_k - x_{k + 1})} = \sum_{k = 1}^{n - 1} \abs{S_{k}} \abs{x_k - x_{k + 1}} \le M \sum_{k = 1}^{n - 1} \abs{x_k - x_{k + 1}}
        $$
        而由题知 $\sum\limits_{n = 1}^{\infty} \abs{x_n - x_{n + 1}}$ 收敛，因此这个和式也收敛。

        于是得证。
    \end{proof}
\end{problem}

\begin{problem}
    习题 11.3.8
    \begin{solution}
        \begin{enumerate}
            \item[\textbf{1)}] 当 $n = 4k$ 时，$\sin \frac{n \pi}{4} = 0$，不考虑。当 $n = 4k + s\ (s \in \{1, 2, 3\})$ 时：
            $$
            \sum_{k = 2}^{\infty} \frac{\sin \frac{(4k + s) \pi}{4}}{\ln n} = \sum_{k = 2}^{\infty} (-1)^{k} \frac{\sin s \pi}{\ln n}
            $$
            这是一个 Lebniz 级数，因此是收敛的。而原级数由 $s = 1, 2, 3$ 这三个收敛级数加括号得到，因此也是收敛的。

            \item[\textbf{2)}]
            \begin{itemize}
                \item 当 $p \le 0$ 时，原级数显然发散。
                \item 当 $0 < p \le 1$ 时，我们有
                $$
                \begin{aligned}
                    \sum_{n = 1}^{N} \cos^2 n & = \sum_{n = 1}^{N} \frac{\cos 2n + 1}{2} \\
                    & = \frac{1}{2} \sum_{n = 1}^{N} \cos 2n + \frac{1}{2} \sum_{n = 1}^{N} (-1)^{n} \\
                    & = \frac{1}{4 \sin 1} \sum_{n = 1}^{N} 2 \sin 1 \cos 2n + \frac{1}{2} \sum_{n = 1}^{N} (-1)^{n} \\
                    & = \frac{1}{4 \sin 1} \sum_{n = 1}^{N} (\sin (2n + 1) - \sin (2n - 1)) + \frac{1}{2} \sum_{n = 1}^{N} \\
                    & = \frac{1}{4 \sin 1} \ab(\sin (2N + 1) - \sin 1) + \frac{1}{2} \sum_{n = 1}^{N} (-1)^{n}
                \end{aligned}
                $$
                因此 $\sum\limits_{n = 1}^{N} \cos^2 n$ 有界，进而 $\sum\limits_{n = 1}^{N} (-1)^{n} \cos^2 n$ 有界。根据 Diriclet 判别法，$\sum\limits_{n = 1}^{\infty} (-1)^{n} \frac{\cos^2 n}{n^{p}}$ 收敛。

                而
                $$
                \begin{aligned}
                    \sum_{n = 1}^{\infty} \abs{(-1)^{n} \frac{\cos 2n}{n^{p}}} & = \sum_{n = 1}^{\infty} \frac{\abs{\cos 2n}}{n^{p}} \\
                    & \ge \sum_{n = 1}^{\infty} \frac{\cos^2 2n}{n^{p}} \\
                    & = \sum_{n = 1}^{\infty} \frac{1 + \cos 4n}{2 n^{p}} \\
                    & \ge \sum_{n = 1}^{\infty} \frac{1}{2 n^{p}} + \sum_{n = 1}^{\infty} \frac{\cos 4n}{2 n^{p}}
                \end{aligned}
                $$
                发散，因此原级数条件收敛。

                \item 当 $p > 1$ 时，有 $\sum\limits_{n = 1}^{\infty} \abs{(-1)^{n} \frac{\cos 2n}{n^{p}}} \le \sum\limits \frac{1}{n^{p}}$，因此原级数绝对收敛。
            \end{itemize}

            \item[\textbf{6)}] 将原式化简：
            $$
            \begin{aligned}
                \sum_{n = 1}^{\infty} \frac{\sin (n + 1) x \cos (n - 1) x}{n^{p}} & = \sum_{n = 1}^{\infty} \frac{\sin 2n x - \sin 2x}{2 n^{p}} \\
                & = \frac{1}{2} \sum_{n = 1}^{\infty} \frac{\sin 2n x}{n^{p}} - \frac{\sin 2x}{2} \sum_{n = 1}^{\infty} \frac{1}{n^{p}}
            \end{aligned}
            $$
            于是我们得到：
            \begin{itemize}
                \item 当 $x = 2k \pi$ 时，原级数收敛，值为 $0$；
                \item 当 $x \neq 2k \pi$ 时：
                \begin{itemize}
                    \item 当 $p \le 1$ 时，根据 \textbf{2)} 的分析可知 $\frac{1}{2} \sum\limits_{n = 1}^{\infty} \frac{\sin 2n x}{n^{p}}$ 收敛，而 $\frac{\sin 2x}{2} \sum\limits_{n = 1}^{\infty} \frac{1}{n^{p}}$ 发散，所以级数发散；
                    \item 当 $p > 1$ 时，有 $\abs{\frac{\sin (n + 1)x \cos (n - 1)x}{n^{p}}} < \frac{1}{n^{p}}$，因此级数绝对收敛。
                \end{itemize}
            \end{itemize}
        \end{enumerate}
    \end{solution}
\end{problem}

\begin{problem}
    习题 11.3.9
    \begin{proof}
        $$
        \sum_{n = 1}^{\infty} \frac{x_n}{n^{a}} = \sum_{n = 1}^{\infty} \frac{x_n}{n^{b}} \cdot \frac{1}{n^{a - b}}
        $$
        因为 $\sum\limits_{n = 1}^{N} \frac{x_n}{n^{b}}$ 有界，$\frac{1}{n^{{a - b}}}$ 单调趋于 $0$，根据 Diriclet 判别法原级数收敛。
    \end{proof}
\end{problem}

\begin{problem}
    习题 11.3.11
    \begin{proof}
        设 $b_n = \ab(1 + \frac{1}{n}) x_n$。使用反证法。假设 $\sum\limits_{n = 1}^{\infty} b_n$ 收敛，那么：
        $$
        \sum_{n = 1}^{\infty} x_n = \sum_{n = 1}^{\infty} \frac{n}{n + 1} b_n
        $$
        因为 $\frac{n}{n + 1}$ 单调有界，$\sum\limits_{n = 1}^{\infty} b_n$ 收敛，根据 Abel 判别法 $\sum\limits_{n = 1}^{\infty} x_n$ 收敛。这与题目矛盾，于是得证。
    \end{proof}
\end{problem}

\begin{problem}
    习题 11.3.12
    \begin{solution}
        \begin{enumerate}
            \item[\textbf{2)}] 考虑
            $$
            \begin{aligned}
                \sum_{n = 1}^{\infty} (-1)^{n} \frac{\cos^2 n}{n} & = \sum_{n = 1}^{\infty} (-1)^{n} \frac{\cos 2n + 1}{2 n} \\
                & = \frac{1}{2} \sum_{n = 1}^{\infty} (-1)^{n} \frac{\cos 2n}{n} + \frac{1}{2} \sum_{n = 1}^{\infty} \frac{(-1)^{n}}{n}
            \end{aligned}
            $$
            结合前面题目的分析可得 $\sum\limits_{n = 1}^{\infty} (-1)^{n} \frac{\cos^2 n}{n}$ 条件收敛。

            同时 $\arctan (3 + n)$ 单调有界，根据 Abel 判别法知原级数条件收敛。

            \item[\textbf{3)}] 考虑
            $$
            \begin{aligned}
                \sum_{n = 2}^{N} (-1)^{n} \sin^2 n & = \sum_{n = 1}^{N} (-1)^{n} \frac{1 - \cos 2n}{2} \\
                & = \frac{1}{2} \sum_{n = 1}^{N} (-1)^{n} - \frac{1}{2} \sum_{n = 1}^{N} (-1)^{n} \cos 2n
            \end{aligned}
            $$
            结合前面题目的分析可得 $\sum\limits_{n = 1}^{N} (-1)^{n} \sin^2 n$ 有界。

            而 $\frac{1}{\ln n}$ 单调趋于 $0$，使用 Diriclet 判别法结合前面题目分析，得到 $\sum\limits_{n = 2}^{\infty} (-1)^{n} \frac{\sin^2 n}{\ln n}$ 条件收敛。

            而 $\ab(1 + \frac{1}{n})^{n},\ (3 - \arctan n)$ 均单调有界，应用两次 Abel 判别法即得到原级数条件收敛。
        \end{enumerate}
    \end{solution}
\end{problem}

\begin{problem}
    习题 12.1.1
    \begin{solution}
        \item[\textbf{1)}] 当 $n \to + \infty$ 时：
        $$
        \frac{n - 1}{n + 1} \ab(\frac{x}{3x + 1})^{n} = O\ab(\ab(\frac{x}{3x + 1})^{n})
        $$
        因此可知该函数项级数收敛当且仅当 $-1 < \frac{x}{3x + 1} < 1$，即收敛域为 $(0, -\frac{1}{2}) \cap (-\frac{1}{4}, +\infty)$。

        \item[\textbf{2)}] 当 $x = 0$ 时，级数收敛。当 $x \neq 0$ 时，$n \to + \infty$ 时有：
        $$
        n (x + \frac{1}{n})^{n} = O(n x^{n})
        $$
        因此该函数项级数收敛域为 $(-1, 1)$。

        \item[\textbf{4)}] 当 $|x| = 1$ 时，级数发散；当 $|x| < 1$ 时，$\frac{x^{n}}{1 + x^{2n}} = O(x^{n})$；当 $|x| > 1$ 时，$\frac{x^{n}}{1 + x^{2n}} = O(\frac{1}{x^{n}})$。因此该函数项级数收敛域为 $\mathbb{R} \setminus \{-1, 1\}$。
    \end{solution}
\end{problem}

\begin{problem}
    习题 12.1.4
    \begin{solution}
        \item[\textbf{2)}] 该函数列逐点收敛于 $0$，考虑 $\sup \abs{S_n(x)}$。
        $$
        S_n'(x) = -2n^2 x^2 (1 - x^2)^{n - 1} + (1 - x^2)^{n} = -n(1 - x^2)^{n - 1} ((1 + 2n) x^2 - 1)
        $$
        因此 $|S_n(x)|$ 在 $[0, 1]$ 上的最大值是 $M_n = S_n\ab(\sqrt{\frac{1}{1 + 2n}}) = n \sqrt{\frac{1}{1 + 2n}} \ab(1 - \frac{1}{1 + 2n})^{n}$。这是收敛的。因此原函数列一致收敛。

        \item[\textbf{4)}] 显然函数列逐点收敛于 $e^{x}$。
        $$
        \begin{aligned}
            \lim_{n \to \infty} \sup_{0 \le x \le a} \abs{e^{x} - S_n(x)} & = \lim_{n \to \infty} \abs{e^{x} - \ab(1 + \frac{x}{n})^{n}} \\
        \end{aligned}
        $$
        而
        $$
        \begin{aligned}
            \ln e^{x} - \ln \ab(1 + \frac{x}{n})^{n} & = x - n \ln \ab(1 + \frac{x}{n}) \\
            & = n \ab(\frac{x}{n} - \ln \ab(1 + \frac{x}{n})) \\
            & \le \frac{x^2}{2n}
        \end{aligned}
        $$
        同时有：
        $$
        \begin{aligned}
            e^{x} - \ab(1 + \frac{x}{n})^{n} & \le e^{\xi} \ab(\ln e^{x} - \ln \ab(1 + \frac{x}{n})^{n}) \\
            & \le e^{a} \cdot \frac{x^2}{2n}
        \end{aligned}
        $$
        因此
        $$
        \lim_{n \to \infty} \sup_{0 \le x \le a} \abs{e^{x} - S_n(x)} \le \lim_{n \to \infty} e^{a} \frac{a^2}{2n} = 0
        $$
        因此原函数列一致收敛。

        \item[\textbf{6)}]
        \begin{itemize}
            \item $x \in (0, +\infty)$：函数列逐点收敛于 $0$，但
            $$
            \sup_{x > 0} \abs{S_n(x) - 0} = \sup_{x > 0} \abs{\frac{1}{1 + n x}} = 1
            $$
            因此该函数列不一致收敛。

            \item $x \in [a, + \infty)$：函数列逐点收敛于 $0$，且
            $$
            \lim_{n \to \infty} \sup_{x > 0} \abs{S_n(x) - 0} = \lim_{n \to \infty} \abs{\frac{1}{1 + a x}} = 0
            $$
            因此该函数列一致收敛。
        \end{itemize}

        \item[\textbf{7)}]
        \begin{itemize}
            \item $x \in (0, 1 - a]$：函数列逐点收敛于 $0$，
            $$
            \lim_{n \to \infty} \sup_{x \in (0, 1 - a]} \abs{\frac{x^{n}}{1 + x^{n}}} = \lim_{n \to \infty} \abs{\frac{(1 - a)^{n}}{1 + (1 - a)^{n}}} = 0
            $$
            因此该函数列一致收敛。
            \item $x \in [1 - a, 1 + a]$：$x = 1$ 时 $\{S_n(1)\}$ 不收敛，因此该函数列不一致收敛。
            \item $x \in [1 + a, +\infty)$：函数列逐点收敛于 $1$。
            $$
            \lim_{n \to \infty} \sup_{x \in [1 + a, +\infty)} \abs{\frac{x^{n}}{1 + x^{n}}} = \lim_{n \to \infty} 1 = 1
            $$
            因此该函数列不一致收敛。
        \end{itemize}

        \item[\textbf{8)}]
        \begin{itemize}
            \item $x \in [-1, 1]$：函数列逐点收敛于 $0$，且
            $$
            \lim_{n \to \infty} \sup_{x \in [-1, 1]} \abs{e^{-(x - n)^2} - 0} = \lim_{n \to \infty} e^{-(1 - n)^2} = 0
            $$
            因此该函数列一致收敛。
            \item $x \in \mathbb{R}$：函数列逐点收敛于 $0$，但 
            $$
            \lim_{n \to \infty} \sup_{x \in \mathbb{R}} \abs{e^{-(x - n)^2} - 0} = 1
            $$
            因此该函数列不一致收敛。
        \end{itemize}
    \end{solution}
\end{problem}

\begin{problem}
    习题 12.1.5
    \begin{solution}
        该函数列逐点收敛于 $0$。
        \begin{enumerate}
            \item[\textbf{1)}]
            $$
            \lim_{n \to \infty} \sup \abs{S_n(x) - 0} = \lim_{n \to \infty} S_{n}\ab(\frac{1}{n}) = \lim_{n \to \infty} \frac{n^{\alpha - 1}}{\E}
            $$
            因此 $\alpha < 1$ 时原函数列一致收敛。

            \item[\textbf{2)}]
            $$
            \begin{gathered}
                \lim_{n \to \infty} \int_0^1 S_n(x) = \lim_{n \to \infty} n^{\alpha} \ab[-\frac{1 + nx}{n^2} \E^{-nx}]_0^1 = \lim_{n \to \infty} n^{\alpha - 2} (1 - (1 + n) e^{-n}) \\
                \int_0^1 \ab(\lim_{n \to \infty} S_n(x)) \dd{x} = \int_0^1 0 \dd{x} = 0
            \end{gathered}
            $$
            因此 $\alpha < 2$ 时两式相等。

            \item[\textbf{3)}]
            $$
            \begin{gathered}
                \lim_{n \to \infty} S'_n(x) = \lim_{n \to \infty} n^{\alpha} (1 - nx) \E^{-nx} \\
                \ab(\lim_{n \to \infty} S_n(x))' = 0
            \end{gathered}
            $$
            两式始终相等。
        \end{enumerate}
    \end{solution}
\end{problem}

\begin{problem}
    习题 12.1.7
    \begin{proof}
        首先因为 $S_0(x)$ 在闭区间连续，其一定有界，记为 $M$。那么可以归纳得到：
        $$
        \abs{S_n(x)} \le \int_a^x M \frac{(x - a)^{n - 1}}{(n - 1)!} \dd{x} \le \frac{M (x - a)^{n}}{n!}
        $$
        因此
        $$
        \lim_{n \to \infty} \sup_{x \in [a, b]} \abs{S_n(x)} = \lim_{n \to \infty} M \frac{(b - a)^{n}}{n!} = 0
        $$
        于是得证。
    \end{proof}
\end{problem}

\begin{problem}
    习题 12.1.9
    \begin{proof}
        由题知
        $$
        \lim_{n \to \infty} S_n(x) = \int_{x}^{x + 1} f(t) \dd{t}
        $$
        而
        $$
        \begin{aligned}
            & \lim_{n \to \infty} \sup_{x \in [a, b]} \abs{S_n(x) - \int_{x}^{x + 1} f(t) \dd{t}} \\
            \le & \lim_{n \to \infty} \sup_{x \in [a, b]} \sum_{i = 1}^{n} \frac{1}{n} \ab(\max_{x \in [a + \frac{i - 1}{n}, a + \frac{i}{n}]} f(x) - \min_{x \in [a + \frac{i - 1}{n}, a + \frac{i}{n}]} f(x))
        \end{aligned}
        $$
        因为 $f(x)$ 在闭区间 $[a, b + 1]$ 上连续，因此也一致连续。所以 $\sup_{x \in [a, b]} (\max \cdot - \min \cdot )$ 趋于 $0$。得证。
    \end{proof}
\end{problem}

\begin{problem}
    习题 12.1.11
    \begin{proof}
        $$
        n \ab(x + \frac{1}{n})^{n} = n x^{n} \ab(1 + \frac{1}{n x})^{n} = n x^{n} \ab(\ab(1 + \frac{1}{n x})^{n x})^{\frac{1}{x}}
        $$
        当 $x \to \pm 1$ 时，上式是无界的。因此函数项级数的一般项不逐点收敛于 $0$，函数项级数不一致收敛。
    \end{proof}
\end{problem}

\begin{problem}
    习题 12.1.13
    \begin{proof}
        该函数项级数逐点收敛于 $0$，但
        $$
        \lim_{n \to \infty} \sup_{x \in (1, +\infty)} \abs{\sum_{i = 1}^{n} \frac{1}{i^{x}}} \ge \lim_{n \to \infty} \abs{n \cdot \frac{1}{n^{\log_n 2}}} = \lim_{n \to \infty} \frac{n}{2} = +\infty
        $$
        因此原函数列不一致收敛。
    \end{proof}
\end{problem}

\begin{problem}
    习题 12.1.14
    \begin{proof}
        对于任意 $n \in \mathbb{N}^{+}$，都有：
        $$
        \begin{aligned}
            \sum_{k = n + 1}^{2n} \frac{\sin k \cdot \frac{\pi}{4 n}}{k} & \ge \sum_{k = n + 1}^{2n} \frac{\sin k \cdot \frac{\pi}{4 k}}{2 n} = \frac{\sqrt{2}}{4} \\
        \end{aligned}
        $$
        这违背柯西收敛原理，因此该函数项级数不一致收敛。
    \end{proof}
\end{problem}
